%========================
% FURRY_Paper_Structure.tex
% TOFE 2026 -- Paper skeleton for assembly in Overleaf.
% Section drafts are written separately and ported in order.
% Status tags: [DRAFTED], [IN PROGRESS], [PENDING]
%========================

\documentclass[]{interact}

\usepackage{graphicx}
\usepackage{epstopdf}
\usepackage{amsmath,amssymb}
\usepackage{booktabs}
\usepackage{multirow}
\usepackage{xcolor}
\usepackage{hyperref}

\usepackage[numbers,sort&compress]{natbib}
\bibpunct[, ]{[}{]}{,}{n}{}{,}

\newcommand{\etal}{\textit{et al.}}
\newcommand{\ie}{i.e.}
\newcommand{\eg}{e.g.}
\newcommand{\INSERT}[1]{\textcolor{red}{\textbf{[#1]}}}

%========================
\begin{document}

\title{Design and Engineering of FURRY: A Plasma-Facing Materials Testing
Platform}

\author{
\name{%
Liam K. King\textsuperscript{a}\thanks{Corresponding author.
Email: lkking3@ncsu.edu}
\and
Florian M. Laggner\textsuperscript{a}
\and
Arthur G. Mazzeo\textsuperscript{a}
\and
Christopher A. Andersen\textsuperscript{a}
\and
Evan Kallenberg\textsuperscript{a,b}
\and
Brendan J. Crowley\textsuperscript{b}
}
\affil{%
\textsuperscript{a}North Carolina State University, Department of Nuclear
Engineering, 2500 Katherine Stinson Dr, Raleigh, NC 27607, USA\\
\textsuperscript{b}General Atomics, P.O.\ Box 85608, San Diego, CA 92186,
USA
}
}

\maketitle

%------------------------
\begin{abstract}
\INSERT{ABSTRACT -- to be drafted last.}
\end{abstract}

\begin{keywords}
beam optics; plasma-facing components; ion extraction; curved grid; IBSimu;
materials testing; space-charge compensation; active cooling
\end{keywords}

\newpage

%========================
\section{Introduction}
\label{sec:introduction}
%========================
% STATUS: [IN PROGRESS]
% Motivation, facility gap (divertor vs. first wall), FURRY concept,
% contributions (three), paper organization.

\INSERT{SECTION 1 -- Introduction. See Sec1\_Introduction.tex}


%========================
\section{Background}
\label{sec:background}
%========================
% STATUS: [DRAFTED]
% CX-neutral physics, first-wall loading (Brooks Fig 1), Tokar CX spectra
% (Fig 2), Eckstein sputtering yield (Fig 3), yield-weighted erosion
% contribution spectrum (Fig 4), FURRY operating rationale.

\INSERT{SECTION 2 -- Background. Port from Core\_draft\_v2.tex \S2.}


%========================
\section{System Overview and Design Objectives}
\label{sec:system}
%========================
% Global-level section. Calculation-light. Establishes what FURRY must
% do and why before discussing how.

\subsection{Mission and Accelerated Exposure Targets}
\label{sec:mission}
% STATUS: [DRAFTED]
% FURRY particle flux, erosion rate, surface recession rate.
% Erosion-equivalent acceleration factor vs. Tokar keV CX reference.
% Three scope limitations stated explicitly.

\INSERT{SECTION 3.1 -- See Sec3p1\_Mission\_Exposure.tex}

\subsection{Operating Envelope}
\label{sec:operating_envelope}
% STATUS: [PENDING]
% Beam energy (up to 10~keV), current (up to 4~A), power (up to 40~kW).
% Average heat flux at target: ~12.6~MW/m^2 at full power.
% Two qualification modes: thermal vs. sputtering-dominant.
% Child--Langmuir perveance scaling for lower-voltage operation.

\INSERT{SECTION 3.2 -- Operating Envelope. PENDING.}

\subsection{Diagnostic Suite}
\label{sec:diagnostics}
% STATUS: [PENDING]
% TDS: hydrogen retention and desorption studies.
% Thermocouples and IR camera: heat flux monitoring.

\INSERT{SECTION 3.3 -- Diagnostic Suite. PENDING.}

\subsection{Steady-State Design Challenges}
\label{sec:design_challenges}
% STATUS: [PENDING]
% Overview of thermal management across all subsystems.
% Ion-optics final design informs downstream subsystem design.
% Motivates feasibility assessments in Section 4.

\INSERT{SECTION 3.4 -- Steady-State Design Challenges. PENDING.}


%========================
\section{Subsystem Engineering Feasibility}
\label{sec:subsystems}
%========================
% Mid-level section. 0D and analytical calculations showing design
% feasibility. Higher-fidelity analysis follows grid design finalization.

\subsection{Ion Source Interface and Faraday Shield}
\label{sec:faraday}
% STATUS: [PENDING]
% LUPIN ICP source overview.
% Faraday shield protecting quartz window under steady-state plasma exposure.
% CFD: 10-channel vs. 30-channel pressure analysis -- 10-channel non-viable
%   (exceeds 60~PSI spec), 30-channel viable at 10~GPM / 24~PSI outlet.
% Thermal: 20~kW surface load -- resin non-viable, copper 30-channel
%   viable with ~20~K temperature variation.

\INSERT{SECTION 4.1 -- Ion Source Interface and Faraday Shield. PENDING.}

\subsection{Sample Holder Active Cooling}
\label{sec:sample_holder}
% STATUS: [PENDING]
% 40~kW onto 63.5~mm diameter target (12.6~MW/m^2).
% Back-fed pin-and-fin copper cooling block.
% 0D thermal feasibility: Q_in = Q_out.
% Pin-fin geometry as heat transfer area multiplier.

\INSERT{SECTION 4.2 -- Sample Holder Active Cooling. PENDING.}

\subsection{Accelerator Grid Thermal Management}
\label{sec:grid_cooling}
% STATUS: [PENDING]
% Screen grid: grounded, full plasma sheath exposure.
%   Transparency ~20% (conservative), plasma power intercepted ~12--16~kW.
%   0D cooling requirement.
% Accelerator grid: ~0.2~A ion interception at 10~kV = ~2~kW deposited.
%   Stated as conservative upper bound.
%   0D cooling requirement.

\INSERT{SECTION 4.3 -- Accelerator Grid Thermal Management. PENDING.}


%========================
\section{Ion-Optics Design}
\label{sec:ion_optics}
%========================
% Primary technical contribution. IBSimu simulations, parameterized scans,
% curved-grid analysis. Highest fidelity content in the paper.

\subsection{Extraction Geometry, Simulation Approach, and Performance Metrics}
\label{sec:methods}
% STATUS: [DRAFTED]
% Flat two-grid and curved-grid geometries, IBSimu model, mesh convergence,
% performance metrics, Child--Langmuir perveance, SCC 0D estimate.

\INSERT{SECTION 5.1 -- Port from Core\_draft\_v2.tex \S3.}

\subsection{Flat Two-Grid Design and Electrostatic Steering Ceiling}
\label{sec:results-steering}
% STATUS: [DRAFTED]
% Parameterized scan results. Max safe steering angle 3.43 deg insufficient
% for full aperture population. Delivered current capped at ~1.39~A.

\INSERT{SECTION 5.2 -- Port from Core\_draft\_v2.tex \S4.1.}

\subsection{Spherically Curved Grid Concept and Geometric Aiming}
\label{sec:results-curved}
% STATUS: [DRAFTED -- INSERT tags pending HPC sweep results]
% Concentric spherical-cap geometry, bore-clearance constraint, central
% beamlet divergence diagnostic. IBSimu sweep over R, r_a, n_i.

\INSERT{SECTION 5.3 -- Port from Core\_draft\_v2.tex \S4.2.}

\subsection{Operating Envelope Characterization}
\label{sec:results-envelope}
% STATUS: [DRAFTED -- INSERT tags pending HPC sweep results]
% Delivered current heatmap in R--n_i space. Recommended design point.
% V--I operating space with constant power contours.

\INSERT{SECTION 5.4 -- Port from Core\_draft\_v2.tex \S4.3.}


%========================
\section{Discussion}
\label{sec:discussion}
%========================
% STATUS: [DRAFTED]
% Performance ceiling set by geometric compatibility.
% Curved-grid as established solution applied to compact test stand.
% Quantitative necessity and sufficiency demonstration.
% Limitations.

\INSERT{SECTION 6 -- Port from Core\_draft\_v2.tex \S5.}


%========================
\section{Conclusions}
\label{sec:conclusions}
%========================
% STATUS: [DRAFTED -- minor revisions needed]
% Flat grid ceiling, curved grid recovery, operating envelope.
% Future work: measured n_i(r), collisional transport, thermal-mechanical
% analysis, prototype fabrication.

\INSERT{SECTION 7 -- Port from Core\_draft\_v2.tex \S6.}


%------------------------
\section*{Acknowledgments}
\INSERT{FUNDING, FACILITY, AND PERSONNEL ACKNOWLEDGMENTS.}

\section*{Author Contributions}
\INSERT{AUTHOR CONTRIBUTION STATEMENT.}

\section*{Data Availability}
\INSERT{DATA AVAILABILITY STATEMENT.}

\section*{Disclosure Statement}
The authors report no conflicts of interest.

%------------------------
\bibliographystyle{tfs}
\bibliography{References}

\end{document}
